Um zu diskutieren, ob $j^0 = \rho$ im Fall der Dirac-Gleichung positiv definit ist, müssen wir die Definition der Wahrscheinlichkeitsdichte $\rho$ im Kontext der Dirac-Gleichung betrachten. Die Dirac-Gleichung beschreibt relativistische Spin-$\frac{1}{2}$-Teilchen und ist gegeben durch:

\[
(i\gamma^\mu \partial_\mu - m)\psi = 0
\]

Hierbei sind $\gamma^\mu$ die Dirac-Matrizen, $\psi$ ist der Dirac-Spinor und $m$ ist die Masse des Teilchens. Die Wahrscheinlichkeitsdichte $\rho$ ist die zeitliche Komponente des Vierer-Stroms $j^\mu$, der durch

\[
j^\mu = \bar{\psi} \gamma^\mu \psi
\]

definiert ist, wobei $\bar{\psi} = \psi^\dagger \gamma^0$ der adjungierte Spinor ist. Insbesondere ist die zeitliche Komponente $j^0$ gegeben durch:

\[
j^0 = \psi^\dagger \psi
\]

Da $\psi^\dagger \psi$ das Skalarprodukt des Spinors $\psi$ mit sich selbst ist, ist $j^0$ stets nicht-negativ, da es sich um die Summe der Betragsquadrate der Komponenten von $\psi$ handelt. Dies impliziert, dass $j^0$ positiv semidefinit ist. 

Für $j^0$ positiv definit zu sein, müsste $j^0 > 0$ für alle $\psi \neq 0$ gelten. Da jedoch $\psi = 0$ möglich ist, ist $j^0$ nicht positiv definit, sondern lediglich positiv semidefinit. 

Zusammenfassend ist $j^0 = \rho$ im Fall der Dirac-Gleichung positiv semidefinit, aber nicht positiv definit.